% Generated draft: requires mathematical and source review.
\documentclass[11pt]{article}
\usepackage[margin=25mm]{geometry}
\usepackage{fontspec}
\setmainfont{Latin Modern Roman}
\usepackage{amsmath,amssymb,mathtools,unicode-math}
\setmathfont{Latin Modern Math}
\usepackage{enumitem,needspace,xurl,hyperref}
\setlength{\parindent}{0pt}
\setlength{\parskip}{5pt}
\setlength{\emergencystretch}{3em}
\newcommand{\N}{\mathbb{N}}
\newcommand{\Z}{\mathbb{Z}}
\newcommand{\Q}{\mathbb{Q}}
\newcommand{\R}{\mathbb{R}}
\newcommand{\C}{\mathbb{C}}
\newcommand{\F}{\mathbb{F}}
\newcommand{\A}{\mathbb{A}}
\newcommand{\PP}{\mathbb{P}}
\newcommand{\E}{\mathbb{E}}
\newcommand{\eps}{\varepsilon}
\newcommand{\dd}{\,\mathrm{d}}
\begin{document}
\section*{UnsolvedMath ID: 30003420}
\textbf{Source problem code:} catalogue entry
\subsection{Definitions and mathematical statement}
Write \(\omega\) for the natural numbers and \(\omega_1\) for the first uncountable ordinal. For \(f,g\in\omega^\omega\), put \(f\leq^*g\) if \(\{n\in\omega:f(n)>g(n)\}\) is finite. The bounding and dominating numbers are
\[
\mathfrak b=\min\{|B|:B\subseteq\omega^\omega\text{ and }(\forall g\in\omega^\omega)(\exists f\in B)\ f\not\leq^*g\},
\]
\[
\mathfrak d=\min\{|D|:D\subseteq\omega^\omega\text{ and }(\forall f\in\omega^\omega)(\exists g\in D)\ f\leq^*g\}.
\]
Let \([\omega]^\omega\) denote the infinite subsets of \(\omega\). A set \(S\subseteq\omega\) splits \(X\in[\omega]^\omega\) if both \(X\cap S\) and \(X\setminus S\) are infinite. Define
\[
\mathfrak s=\min\{|\mathcal S|:\mathcal S\subseteq[\omega]^\omega\text{ and }(\forall X\in[\omega]^\omega)(\exists S\in\mathcal S)\ S\text{ splits }X\}.
\]
A family \(\mathcal A\subseteq[\omega]^\omega\) is almost disjoint if distinct members have finite intersection. It is maximal almost disjoint if every \(X\in[\omega]^\omega\) has infinite intersection with some member of \(\mathcal A\). Define
\[
\mathfrak a=\min\{|\mathcal A|:\mathcal A\text{ is an infinite maximal almost disjoint family on }\omega\}.
\]
The requirement that \(\mathcal A\) be infinite excludes finite maximal almost disjoint families.

For an infinite cardinal \(\kappa\), a proper ultrafilter \(U\) on \(\kappa\) is uniform if every member of \(U\) has cardinality \(\kappa\). A base for \(U\) is a family \(\mathcal B\subseteq U\) such that
\[
(\forall A\in U)(\exists B\in\mathcal B)\ B\subseteq A.
\]
Under the uniform-ultrafilter convention, define
\[
\mathfrak u(\kappa)=\min\{|\mathcal B|:\mathcal B\text{ is a base for a uniform ultrafilter on }\kappa\}.
\]
Here \(2^{\aleph_1}=|\mathcal P(\omega_1)|\).

Using these standard interpretations, the supplied record asks the following three consistency questions, with all characteristics computed inside the model in question:
\begin{enumerate}
\item Can a model of \(\mathrm{ZFC}\) satisfy \(\mathfrak b<\mathfrak s<\mathfrak d\)?
\item Can a model of \(\mathrm{ZFC}\) satisfy \(\mathfrak b<\mathfrak a<\mathfrak s\)?
\item Can a model of \(\mathrm{ZFC}\) satisfy \(\mathfrak u(\omega_1)<2^{\aleph_1}\)?
\end{enumerate}
The record supplies no definitions of these characteristics. In particular, the third question above is conditional on the intended meaning of \(\mathfrak u(\omega_1)\) being the uniform-ultrafilter number with the stated base convention. The record also does not specify a relative-consistency hypothesis, permissible additional axioms, or large-cardinal assumptions. Taking \(\mathrm{ZFC}\) as the ambient theory is an explicit conventional interpretation, not a supplied consistency-strength assertion. The questions are presented separately; the record does not explicitly ask for one model satisfying all three inequalities simultaneously.
\subsection{Short English statement}
The record asks whether three strict inequality patterns among cardinal characteristics can occur in models of set theory. The first two concern bounding, splitting, dominating, and almost disjointness on \(\omega\); the third concerns the size of an ultrafilter base on \(\omega_1\), interpreted here using uniform ultrafilters. The exact higher-cardinal convention and consistency assumptions require source confirmation. The supplied status is partially solved; this draft does not independently establish that assessment.
\subsection{Sources}
\begin{itemize}
\item Supplied local registry record for permanent UnsolvedMath ID \(30003420\), source problem code \texttt{catalogue entry}, in the project's \texttt{data/unsolvedmath/problems.json} and supplementary \texttt{catalog/problems.json} registry system. Title: \emph{Consistency Patterns Among Cardinal Characteristics}. Collection: Oberwolfach Reports Open Problems. No individual proposer is supplied; the proposed year is recorded as 2017.
\item Supplied report citation: \emph{Set Theory (2018)}, Oberwolfach Reports, DOI \texttt{10.4171/owr/2017/11}, pp.~527--589. \url{https://doi.org/10.4171/owr/2017/11}. The supplied title year differs from the year appearing in the DOI and proposed-year field; this discrepancy is retained pending verification.
\item Supplied literature source: \emph{Cardinal characteristics at aleph omega}. \url{https://arxiv.org/abs/1812.05291}.
\item Supplied literature source: \emph{OWR Set Theory report}. \url{https://ems.press/content/serial-article-files/46674}.
\item Literature assessment dated exactly \texttt{2026-08-21}: supplied dated curation, not independently verified. The record labels the problem \texttt{partially\_solved} and reports related forcing separations, including small ultrafilter numbers at higher cardinals, while stating that the full list of three extracted consistency questions is not settled as one theorem. This assessment does not identify which individual questions are resolved or specify their consistency hypotheses. The background expressly calls for independent verification of time-sensitive status and cited claims.
\item The supplied statement status is \texttt{corrected\_verified}; its accompanying explanation says that mathematical content was preserved while notation, typography, and sentence structure were repaired for publication. This is a supplied editorial qualification, not independent mathematical verification of this draft.
\end{itemize}
\end{document}
